\documentclass{ximera}


\begin{document}
\author{Alexander Holvoet}
\xmtitle{Een ronde schijf en een vierkant als grondvlak}


\subsection*{Inleiding}

Hier komen de lesactiviteiten van UW4201 (zie comments) als Ximera.
Open vraag: hoe integreren we de leerkrachtenadviezen? of gwn verwijzen naar magazine dan?
Open vraag: hoe integreren we de bronnen?

% \begin{lesactiviteit}{Een hoop zand op een ronde ondergrond}

% Giet zo veel mogelijk zand of zout op een ronde basisvorm, bijvoorbeeld een rond bierviltje of een omgekeerd plastieken bekertje. Je moet het niet al te voorzichtig doen: wat te veel is, glijdt er af. Daarom staat je basisvorm best een beetje hoger.
% \newpage
    
% 	\begin{vraagenantwoord}
%     	\vraag{Welke vorm hebben de hopen?}
%         \antwoord{Het zijn allemaal kegels. Je zult merken dat deze kegels niet noodzakelijk perfect zijn: de top is niet juist één punt. Hier kun je uitleggen dat we werken met een \textit{wiskundig model} dat vrij goed aansluit bij de realiteit maar er niet mee samenvalt.}
%     \end{vraagenantwoord}

   

    
%  We definiëren een nieuw begrip: de \emph{rusthoek} is de hoek tussen de beschrijvende rechten van de kegel en de horizontale richting. Die hoek is over de hele flank gelijk, dus niet enkel vanonder, zie de figuur hieronder.
%     \afbeelding[5cm]{UW4201/img/Lrusthoek}{}{} 


% \begin{vraagenantwoord}[resume]
		
% 	\vraag{Denk even na hoe je deze hoek zo nauwkeurig mogelijk kunt meten.}
%     \antwoord{ Uiteraard kun je een geodriehoek gebruiken. Dit is niet zo eenvoudig als het op het eerste zicht lijkt. Zodra je de hoop raakt, begint die te schuiven... Je zou ook een foto kunnen nemen van de kegel in vooraanzicht en daarop de rusthoek meten of je kunt de foto in GeoGebra plakken en daarop de rusthoek meten. Of je kunt de straal van het grondvlak en de hoogte van de kegel meten en hiermee de tangens van de rusthoek en dus ook de rusthoek zelf berekenen.}
%      \afbeelding[6cm]{UW4201/img/rusthoek meten met geogebra}{}{} 
%     \vraag{Meet en noteer de rusthoeken voor hopen met dezelfde diameter maar met verschillend materiaal. Verdeel eventueel het werk.}
%     \antwoord{Bijvoorbeeld:}
%         \begin{tabel}
% 		\begin{tabular}{m{3cm} m{3cm} m{2cm}}  
% 			\tableheader{diameter} & \tableheader{materiaal} & \tableheader{rusthoek}  \\ 
% 		10 cm & zand   & 32° \ \\
% 		10 cm & rijst  & 25° \\
% 		10 cm  & zout & 30° \\
%     	\end{tabular}
% 	\end{tabel}
    
%     \afbeelding[6cm]{UW4201/img/Lmateriaal.jpg}{}{} 
    
%     \vraag{Formuleer een besluit.}
% 	\antwoord{Als je de diameter van het grondvlak constant houdt, dan hangt de rusthoek af van het soort materiaal.}
    
%     \vraag{Doe nu hetzelfde voor hopen van hetzelfde materiaal maar met verschillende diameters. Meet en noteer ook hier je meetresultaten in een tabel. Wat besluit je?}
   
%        \antwoord{Bijvoorbeeld:
%        \begin{tabel}
% 		\begin{tabular}{m{3cm} m{3cm} m{2cm} }  % todo: moet eigenlijk 7.5 cm breed zijn
% 			\tableheader{materiaal} & \tableheader{diameter} & \tableheader{rusthoek} \\
% 			zand&10 cm    & 32°  \\
% 			zand&15 cm  & 32° \\
% 			zand&20 cm   & 32° \\
%     	\end{tabular}
% 	\end{tabel}
%     De rusthoek hangt niet af van de diameter van het grondvlak.}
% 	\end{vraagenantwoord}
	
% \end{lesactiviteit}

% \begin{lesactiviteit}{Een hoop zand op een vierkant}
	
% 	\begin{vraagenantwoord}
%     \vraag{Wat denk je te krijgen als je zand giet op een vierkant plaatje? Doe een voorspelling en voer daarna pas het experiment uit.}
%    \antwoord{Je krijgt een rechte piramide met het vierkant als grondvlak. }
%    \afbeelding [5cm]{UW4201/img/LPiramide}{}{fig:piramide} 

% \vraag{Meet opnieuw de rusthoek en vergelijk met de rusthoek bij de kegel.}
% \antwoord{De rusthoek is dezelfde  als bij de kegel, tenminste als je hetzelfde materiaal gebruikt hebt.}
% \vraag{Teken het bovenaanzicht van de zandhopen met een rond grondvlak en een vierkant grondvlak.}
%   \antwoord{Bij het ronde grondvlak krijg je een cirkel met een puntje. Bij het vierkante grondvlak krijg je een vierkant met diagonalen. Het puntje staat voor de top van de kegel, de diagonalen staan voor de ribben van de piramide. Uit het bovenaanzicht kun je niet afleiden hoe groot de rusthoek is.}    
%   \afbeelding[7cm]{UW4201/img/Lbovenzichtkegel}{}{fig:bovenzicht}
%  \end{vraagenantwoord}
% \end{lesactiviteit} 

% \begin{lesactiviteit}{Een hoop zand op een rechthoekig grondvlak}

%   \begin{vraagenantwoord}
%       \vraag{Welke zandhoop verwacht je bij een rechthoekig grondvlak? Maak een schets van het te verwachten bovenaanzicht en voer daarna pas het experiment uit.}
%       \antwoord{Je krijgt geen piramide meer. Er is niet één top maar een lijnstuk met hoogste punten. (In de bouw kom je dit soort ruimtelijk lichaam tegen in de vorm van een \textit{schilddak}.) }
%       \vraag{Wat kun je zeggen over de rusthoek?}
%       \antwoord{Deze is weer constant en gelijk aan die  bij de kegel en de piramide met vierkant grondvlak. De rusthoek hangt immers enkel af van het soort materiaal.}
%       \end{vraagenantwoord}

% Hoogtelijnen ken je uit de aardrijkskundeles: het zijn lijnen op een kaart die alle punten met een gelijke hoogte boven het zeeniveau verbinden. 
%       \begin{vraagenantwoord}[resume]
      
%       \vraag{Vul je tekening van het bovenaanzicht aan met enkele hoogtelijnen.}
%       \antwoord{De hoogtelijnen zijn evenwijdig met de zijden van de rechthoek en de afstand tussen twee opeenvolgende hoogtelijnen  is altijd dezelfde. Dit komt omdat de flanken overal even steil zijn. }
             
%       \afbeelding[6cm]{UW4201/img/LRHhoogtelijnen.png}{}{}
%  \end{vraagenantwoord}
%       De hoogtelijn met de hoogste punten noemen we net zoals bij echte bergen de kamlijn. Ze is evenwijdig met de langste zijde van het grondvlak en staat loodrecht op de kortste zijde. Je ziet dat die in het midden ligt. Het kan zijn dat je die niet getekend hebt. 
% Verderop zullen we proberen de positie van de kamlijn meetkundig te bepalen.

% \begin{vraagenantwoord}[resume]


% \vraag{Verbind de hoekpunten van de hoogtelijnen.} \antwoord{Je krijgt volgende figuur.
%     \afbeelding[6cm]{UW4201/img/LRHhoogtelijnenbiss.png}{}{}}

   
%  \clearpage
% \vraag{Aan welke eigenschap voldoen a) punten die op de kamlijn liggen zoals het punt $R$; b) punten op één van de schuine rechten die de hoekpunten van de rechthoekige hoogtelijnen verbinden zoals het punt $P$;  c) punten die op zowel de kamlijn als op de twee schuine rechten liggen zoals het punt $S$? }
%    \afbeelding[6cm]{UW4201/img/LRHbiss.png}{}{}
% \antwoord{a) Punten op de kamlijn bevinden zich even ver van de twee overstaande zijden, zowel in 3D als in het bovenaanzicht. Dit verklaart waarom de kamlijn in het bovenaanzicht in het midden tussen de twee lange randen ligt. b) Deze punten liggen even ver van twee opeenvolgende randen, ook weer zowel in 3D als in bovenaanzicht. Dit verklaart waarom deze lijnen een hoek van 45° maken met deze zijden. c) Deze twee punten bevinden zich even ver van drie opeenvolgende randen.}
% \end{vraagenantwoord}
% In het bovenaanzicht is de \textit{kamlijn} dus de meetkundige plaats van de punten die zich even ver bevinden van de twee dichtstbijzijnde zijden. De punten op het uiteinde van de kamlijn liggen op gelijke afstand van drie zijden. In 3D kunnen we de ribben daarom ook kamlijnen noemen.  
% \end {lesactiviteit}

% \begin{lesactiviteit}{Een hoop zand op een driehoekig grondvlak}

%   \begin{vraagenantwoord}
     
%       \vraag{Neem als volgende grondvlak een willekeurige driehoek. Schets eerst wat je verwacht als bovenaanzicht te verkrijgen, teken ook de vermoedelijke hoogtelijnen. Giet dan zoveel mogelijk zand op de omhoog geplaatste driehoek. Komt het experiment overeen met je schets? Is de rusthoek ook hier weer gelijk aan die in de vorige experimenten? }
%       \antwoord{Je krijgt een piramide, zoals je in de volgende foto's ziet. Die is enkel afhankelijk van de vorm van het grondvlak en niet van waar je begint met zand te gieten. De rusthoek is weer dezelfde.}
%       \begin{afbeeldingenv}% workaround vanuit preambule om foto's naast elkaar te krijgen
% 		\includegraphics[width=6cm,keepaspectratio]{UW4201/img/Lzandopdriehoek.JPG}
%         \quad
%         \includegraphics[width=6cm,keepaspectratio]{UW4201/img/Lhoogtelijnendriehoek.png}
% 	\end{afbeeldingenv}
%   \vraag{Als je ook hier de hoekpunten van de hoogtelijnen verbindt, zie je speciale lijnen verschijnen op het bovenaanzicht. Welke lijnen zijn dit? }   
%   \antwoord{Het zijn de bissectrices van de hoeken gevormd door twee opeenvolgende zijden van de driehoek.
%     \afbeelding[6cm]{UW4201/img/Ldriehoekkortsteweg.png}{}{}}
%  \vraag{Geef hier een verklaring voor. Waarmee komen deze lijnen overeen in 3D?}
%  \antwoord{De bissectrices zijn de meetkundige plaats van de punten die even ver liggen van twee opeenvolgende zijden van de driehoek. In 3D komen ze overeen met de ribben van de piramide. Korrels op de ribben van de piramide weten niet naar welke kant ze moeten vallen en blijven dus liggen.}
%  \vraag{We weten dat in een vlak de drie bissectrices van een driehoek elkaar snijden in één punt. Wat weet je van dit snijpunt?}
%  \antwoord{Omdat het snijpunt op elk van de bissectrices ligt, ligt het even ver van de drie zijden van de driehoek. }
%  \end{vraagenantwoord}
% \clearcolumn
% De bissectrices van het bovenaanzicht stellen in de zandhoop alle punten voor op de ribben van de piramide. Deze punten bevinden zich even ver van twee randen onderaan de zandhoop. De loodlijnen op de zijden geven de kortste weg aan voor een zandkorrel om naar beneden te rollen.
% \afbeelding[8cm]{UW4201/img/zandopdriehoek.png}{}{}
% \begin{vraagenantwoord}[resume]

% \vraag{Wat betekent het snijpunt van de bissectrices uit het bovenaanzicht voor de zandhoop?}
% \antwoord{Dit snijpunt is de projectie van de top van de zandhoop. Deze top is het punt dat zich even ver van de drie randen van het grondvlak bevindt.}

% \end{vraagenantwoord}
% \end{lesactiviteit}

% %\begin{lesactiviteit}{Convexe grondvlakken}
% Het volgende stuk is niet helemaal uitgeschreven als lesactiviteit maar het is wel de bedoeling om dit op een onderzoekende manier aan te brengen te zoals in de vorige paragrafen. Dit kan in groepjes of via een onderwijsleergesprek omdat het wat moeilijker wordt. Het interessante van deze situatie is niet zozeer het eindresultaat maar vooral de zoektocht ernaartoe. 

% We bekijken enkele trapezia. Een trapezium is een vierhoek met twee evenwijdige zijden. We vertrekken van een driehoek. Teken de bissectrices en de ingeschreven cirkel. Doe dit in viervoud.

% Knip nu van drie driehoeken een stuk volgens een lijn evenwijdig met de basis van de driehoek zoals in figuur 3 is aangegeven: boven de ingeschreven cirkel (situatie 2), rakend aan de ingeschreven cirkel (situatie 3) en dwars door de ingeschreven cirkel (situatie 4).
% \afbeelding [15cm]{UW4201/img/Ltrapezia.png}{Een driehoek en drie verschillende soorten trapezia}{}

% Teken nu ook nog de bissectrices van de twee nieuw verschenen hoeken van de trapezia.

% Voor je het experiment uitvoert, kun je uiteraard proberen te voorspellen welke zandhoop en welk bovenaanzicht je zult krijgen. Het loont echt de moeite om dit zelf verder uit te zoeken.

% Als je zand giet op de volledige driehoek (situatie 1), zul je uiteraard een piramide krijgen. We bekijken ook situatie 2 en 4 in detail; situatie 3 laten we aan de lezer.

% In de tweede situatie zijn de bissectrices niet concurrent. Er zijn twee snijpunten van telkens twee bissectrices, de punten $O$ en $O'$, zie figuur 4. De punten $O$ en $O'$ liggen op de bissectrice van de afgesneden hoek van de driehoek. Deze punten zijn de projecties van $S$ en $S’$, de eindpunten van de kamlijn. 
% \afbeelding[7cm]{UW4201/img/Ltrap2.png}{Projectie van situatie 2}{}
% \afbeelding[9cm]{UW4201/img/Ltrap23D.png}{De 3D-versie van situatie 2}{}

% In de 3D-tekening zien we twee rechthoekige driehoeken $HSO$ en $H’S’O’$ die eenzelfde rusthoek hebben en dus gelijkvormig zijn. We zien dat $|OS| = |OH|\cdot $tan$\theta$ en $|O’S’| = |O’H’|\cdot $tan$\theta$ maar $|OH| > |O’H’|$, dus $|OS|>|O’S’|$. 

% Het lijnstuk [SS’] loopt dus niet horizontaal. In het bovenaanzicht kun je dit zien aan het feit dat $O'$ dichter bij de rand ligt dan $O$. Verifieer het met het experiment. 

% En meer zelfs: $|OS|-|O’S’| = (|OH|-|O’H’|)\cdot$ tan$\theta$.

% Als we de kniplijn in onze oorspronkelijke driehoek helemaal zouden opschuiven tot in het hoekpunt, dan zou $|O’H’|$ gelijk zijn aan 0, dus de helling van $[SS’]$ zou dezelfde zijn als die van $[SA]$. Als we $[SS']$ verlengen, zou die door de top $A$ van de driehoek lopen die we hebben afgeknipt.

% \afbeelding[7cm]{UW4201/img/Lsituatie2.jpg}{Situatie 2}{}


% In situatie 4 hebben we ook de vier bissectrices van de hoeken van het trapezium getekend. Nu snijdt telkens één van de oude bissectrices met één van de nieuwe bissectrices. Dit geeft weer twee punten O en O'. Het lijnstuk $[OO']$ dat deze twee punten verbindt, is evenwijdig met de twee evenwijdige zijden van het trapezium. $O$ en $O’$ zijn de projecties op het grondvlak van de toppen $S$ en $S’$.
% \afbeelding[6cm]{UW4201/img/Ltrap4.png}{Projectie van situatie 4}{}
% \afbeelding[7cm]{UW4201/img/Ltrap43D.png}{3D-tekening van situatie 4}{}
% In de 3D-voorstelling zijn er weer twee driehoeken geplaatst, $HSO$ en $H’S’O’$, die twee keer dezelfde rusthoek hebben met het grondvlak. We vinden:
% $|OS| = |OH|\cdot \tan{\theta}$ en $|O’S’| = |O’H’|\cdot \tan{\theta}$ maar nu is $|OH| =|O’H’|$ en dus $|OS| = |O’S’|$.

% Dus $[SS’]  \parallel [HH’] \parallel [KK’]$.
% De kamlijn is evenwijdig met de evenwijdige zijden van het trapezium. Verifieer weer met het experiment.
% %\end{lesactiviteit}

% \begin{lesactiviteit}{Een oefening}
% \begin{vraagenantwoord}

% \vraag{Construeer door gebruik te maken van bissectrices de bovenaanzichten van de zandhopen met volgende grondvlakken. Je kunt hiervoor al of niet GeoGebra gebruiken. Verifieer daarna met een experiment.}
% \afbeelding[14cm]{UW4201/img/Lconvexeveelhoeken.png}{}{}
%  \antwoord{Zie hieronder de oplossingen. De aardrijkskundige hoogtelijnen uit de 3D-figuren zijn erbij getekend, maar ze zijn niet noodzakelijk om de oefening op te lossen.
% \afbeelding[14cm]{UW4201/img/Lconvexemethoogtel.png}{}{}  
%     }
% \end{vraagenantwoord} 
% \end{lesactiviteit}

% \begin{lesactiviteit}{Zand op een schijf met een gat in}

% Het grondvlak waarop we zand willen gieten, is een cirkelschijf met een gaatje in.

% Je ziet hieronder twee cirkelschijven: één met een gat in het midden, één met een gedecentreerd gat. Het gat is klein en in de redeneringen zullen we het als een punt behandelen. Maar in werkelijkheid mag het gat ook niet \emph{te} klein zijn, want de zandkorrels moeten erdoor kunnen.

% \begin{multicols}{2}

% \afbeelding[4cm]{UW4201/img/Zand 07.JPG}{}{}

% \afbeelding[4cm]{UW4201/img/Zand 08.JPG}{}{}

% \end{multicols}

% \begin{vraagenantwoord}
   
%     \vraag{Teken op deze bovenaanzichten telkens de kamlijn.}

%     \vraag{Giet dan zand op het kartonnen exemplaar om dit te verifiëren.}

%     \vraag{Verklaar ten slotte de vorm van de kamlijn.}
    

%     \antwoord{Bij het centrale gat is het niet moeilijk: het bovenaanzicht van de kamlijn is een concentrische cirkel met als straal de helft van de straal van de gegeven schijf. Immers, deze cirkel bestaat uit de punten op gelijke afstand tot het gat en de rand. Dit is ook in 3D te verklaren: in het midden van de kegel heb je een kegelvormige uitsparing met dezelfde rusthoek. Hieronder zie je het bovenaanzicht en het vooraanzicht, en ernaast de kegel met een tweede, gelijke kegel maar op zijn kop (of liever: op zijn top). Deze tweede kegel bepaalt de uitsparing in de eerste kegel.} 

% \begin{multicols}{2}
    
   
% \afbeelding[4cm]{UW4201/img/Zand 09.JPG}{}{}

% $$  $$



% \afbeelding[6cm]{UW4201/img/Zand 10.JPG}{}{}



% \end{multicols}

% \emph{Bij de cirkelschijf met een gedecentreerd gat is de kamlijn in bovenaanzicht een ellips. Als het gat dicht bij het middelpunt ligt, zou je op het eerste gezicht nog kunnen denken dat het een cirkel is. }

% \emph{Maar laten we redeneren: een willekeurig punt $P$ van de kamlijn ligt even ver van het gat $G$ als van de rand van de schijf. Die afstand $|PQ|$ tot de rand van de schijf (zie figuur hieronder) wordt loodrecht op die rand gemeten, dus op een straal van de cirkel. Bijgevolg is 
% $$|PM|+|PG|=|PM|+|PQ|=|MQ|=r$$
% waarbij $r$ de straal is van de cirkel. }

% \afbeelding[4cm]{UW4201/img/Zand 11.JPG}{}{}

% \emph{Dus is de som $|PM|+|PG|$ constant. De punten $P$ van de kamlijn vormen dus een ellips met als brandpunten $G$ en $M$.}


% \vraag{Je kunt, bij het gedecentreerde gat, ook redeneren met hoogtelijnen. Je ziet hier de cirkelschijf met het gat, en concentrische cirkels rond het middelpunt en rond het gat. Construeer met deze figuur de kamlijn en leg uit.}

% \afbeelding[4cm]{UW4201/img/Zand 12.JPG}{}{}

% \antwoord{Je kunt in het vlak redeneren: de punten die we aanduiden moeten op gelijke afstand liggen van het gat als van de rand van de cirkel. Bijvoorbeeld: 5 eenheden van het gat en 5 eenheden van de rand, waarbij deze eenheden de afstanden zijn tussen de concentrische cirkels. Dit geeft de punten $P$ en $P'$. Analoog voor andere eenhedenaantallen.}


% \afbeelding[4cm]{UW4201/img/Zand 13.JPG}{}{}


% \emph{Je kunt ook in 3D redeneren: de cirkels zijn hoogtelijnen en je snijdt de cirkels die even hoog zijn. De cirkels rond het gat beginnen in het gat met hoogte nul (cf. de kegel op zijn top); de cirkels rond het middelpunt beginnen aan de rand met hoogte nul (cf. de gewone kegel).} 



% %\antwoord{Op de figuur hieronder: de zandhoop is een kegel met de vaste rusthoek, waaruit een tweede kegel, omgekeerd en met dezelfde rusthoek, is 'weggehaald'. De top van de weggehaalde kegel is het gat. De 3D-versie van de kamlijn is waar deze twee kegels elkaar snijden. ??Kun je met deze 3D-figuur ook uitleggen dat de doorsnede vlak is?? }

% %\afbeelding[6cm]{UW4201/img/Zand ksn 6 bis.JPG}{}{}

% \vraag{Het gat was in de vorige opgaven heel klein en we behandelden het als een punt. Wat verandert er als het (gedecentreerde) gat een grotere cirkel(schijf) is?} 

% \antwoord{Het bovenaanzicht van de kamlijn is opnieuw een ellips. De brandpunten zijn het middelpunt van de cirkel en het middelpunt van het gat. De som van de afstanden tot deze twee punten is weer constant. Deze constante is niet meer de straal van de cirkel, maar de som van de straal van de cirkel en die van het gat. }

% \emph{Laten we dit even aantonen. Een punt $P$ van de kamlijn (in bovenaanzicht) ligt even ver van de rand van de grote cirkel als van de rand van het ronde gat. Met de notaties van de figuur hieronder: |PQ|=|PH|. De som van de afstanden van $P$ tot het middelpunt $M$ en het middelpunt $G$ van het gat is bijgevolg:}
% $$|PM|+|PG|=|PM|+|PH|+|HG|=|PM|+|PQ|+|HG|=|MQ|+|HG|=r+s$$
% \emph{waarbij $r$ de straal is van de grote cirkel en $s$ de straal van het gat. Deze som is dus constant en $P$ beschrijft een ellips.}
% \begin{multicols}{2}
    

% \afbeelding[4cm]{UW4201/img/Zand 14.JPG}{}{}

% \afbeelding[6cm]{UW4201/img/Zand 15.jpg}{}{} 

% \end{multicols}

% \end{vraagenantwoord}

% \end{lesactiviteit}

% \begin{lesactiviteit}{Zand op een cirkelsegment}

% \afbeelding[5cm]{UW4201/img/Zand 16.JPG}{}{}

% We nemen hier als basis voor een zandhoop een deel van een cirkelschijf aan één kant van een koorde. Deze figuur noemen we een \emph{cirkelsegment}. Het weggesneden stuk is ook een cirkelsegment. Wij nemen hier een `groot' cirkelsegment $-$ dit wil zeggen het stuk met het middelpunt erop $-$ maar het gaat even goed met een `klein' cirkelsegment.

% \begin{vraagenantwoord}

%     \vraag{Teken  (het bovenaanzicht van) de kamlijn die je verwacht.}
    
%     \vraag{Giet dan zand op het kartonnen exemplaar om dit te verifiëren.}

%     \vraag{Waar voldoen de punten van de kamlijn aan? }

%     \antwoord{Voor elk punt $P$ van de kamlijn is de afstand $|PQ|$ tot de koorde  gelijk aan de afstand $|PR|$ tot de rand van de cirkel (zie figuur hieronder).}

%    \afbeelding[5cm]{UW4201/img/Zand 17.JPG}{}{}
%     \end{vraagenantwoord}

% Je vermoedt waarschijnlijk dat (het bovenaanzicht van) de kamlijn een parabool is. In de volgende opgaven ga je dit bewijzen.

%     \begin{vraagenantwoord}[resume]

%     \vraag{Om te bewijzen dat de kamlijn een parabool is, heb je een brandpunt en een richtlijn nodig. Het brandpunt zal het middelpunt $M$ van de cirkel zijn. Waar ongeveer moet de richtlijn $d$ liggen? Toon dit met je vinger.}

%     \antwoord{Leerlingen tonen een rechte evenwijdig met de koorde (vanwege de symmetrie) en 'onder' de kamlijn. De precieze plaats van de richtlijn komt in de volgende vraag aan bod.}
    
%  \vraag{Bewijs door een geschikte keuze van de richtlijn dat de kamlijn inderdaad een parabool is.}

%     \antwoord{Neem een willekeurig punt $P$ van de kamlijn en noem $r$ de straal van de cirkel. We willen de richtlijn $d$ zo bepalen, dat de afstand $|PM|$ tot het brandpunt gelijk is aan de afstand $|PS|$ tot de rechte $d$. We weten dat $|PR|=|PQ|$ (zie vraag 3) en dus:}
%    $$|PM|=r-|PR|=r-|PQ|.$$
%     \afbeelding[5cm]{UW4201/img/Zand 18.JPG}{}{}
    
% \emph{Dus, om te verkrijgen dat $|PM|= |PS|$ (met de benamingen van de figuur) moeten we $d$ op afstand $r$ van de koorde nemen. En dan geldt inderdaad}
%   $$|PM|=r-|PR|=r-|PQ|=|QS|- |PQ|=|PS|.$$
%     %\antwoord{De kamlijn is een parabool als elk punt van de kamlijn even ver ligt van $M$ als van $d$. De leerlingen kunnen bijvoorbeeld de top gebruiken om te bepalen waar de rechte $d$ moet liggen (figuur hieronder links). }

%  %\begin{multicols}{2}

% % \afbeelding[4cm]{UW4201/img/Zand ksn 10bis}{}{} 
 
% % \afbeelding[4cm]{UW4201/img/Zand ksn 8bis}{}{}

%     %   \end{multicols} 

%    % \emph{De afstand tussen deze rechte $d$ en de koorde is gelijk aan de straal van de cirkel. Immers, de top behoort tot de kamlijn en ligt dus even ver van de koorde als van de cirkel.}

% %\vraag{Bewijs nu dat de kamlijn een parabool is.}

%     %\antwoord{We moeten bewijzen dat $|PM|=d(P,d)$. De afstand $|PM|$ is de straal $|RM|$ min de afstand $|PR|$ tot de rand van de cirkel: }

% %$$|PM|=|RM|-|PR|=r-|PR|$$
% %Maar omdat $P$ op de kamlijn ligt, is $|PR|=|PQ|$ en hebben we
% %$$|PM|=r-|PQ|$$
% %De richtlijn $d$ moet zo zijn dat 
% %$$d(P,d)=|PS|=|PM|=$$

% \vraag{Construeer de kamlijn met hoogtelijnen.}

% \afbeelding[5cm]{UW4201/img/Zand 19}{}{}

% \antwoord{Zie hieronder. Op de eerste figuur worden de snijpunten van de overeenkomstige hoogtelijnen geconstrueerd. De rechte hoogtelijnen horen bij een schuin vlak dat volgens de rusthoek afloopt naar de koorde van het grondvlak. De cirkelvormige hoogtelijnen horen bij de kegel die je zou hebben als de basis de volledige cirkelschijf zou zijn. De kamlijn is de doorsnede van de kegel met dat schuine vlak (even schuin als de beschrijvenden van de kegel). Dit is nog een extra manier om in te zien dat het een stuk parabool is. De tweede figuur bevat de uiteindelijke hoogtelijnen van de zandhoop.}

% \begin{multicols}{3}
    
%   \afbeelding[5cm]{UW4201/img/Zand 20.JPG}{}{}

%     \afbeelding[5cm]{UW4201/img/Zand 21.JPG}{}{}

%     \afbeelding[5cm]{UW4201/img/Zand 22}{}{}
        
%     \end{multicols}
    
% \end{vraagenantwoord}
% \end{lesactiviteit}

% \begin{lesactiviteit}{Zand op een niet-convexe veelhoek}

% \begin{vraagenantwoord}

%     \vraag{Teken de kamlijnen die je verwacht als je op deze figuur zand zou gieten.}

%     \afbeelding[7cm]{UW4201/img/Zand 26.JPG}{}{}

%     \antwoord{We verwachten dat veel leerlingen de volgende lijnen zullen tekenen.}

%     \begin{multicols}{2}
%         \afbeelding[7cm]{UW4201/img/Zand 27.JPG}{}{}
%         \afbeelding[7cm]{UW4201/img/Zand 30.JPG}{}{}
%     \end{multicols}

%     \emph{Als ze dit met zand verifiëren, zien ze dat de kamlijn niet alleen uit rechte lijnstukken bestaat. In zijaanzicht zien ze dat de zandhoop in de bocht ook wat hoger komt. Dit brengt hen ertoe om er grondiger over na te denken.}

%     \begin{multicols}{2}
    
%      \afbeelding[5cm]{UW4201/img/Zand 28.JPG}{}{}
%      \afbeelding{UW4201/img/Zand 29.JPG}{}{}
        
%     \end{multicols}

%     \emph{De stukken van de kamlijn kunnen niet allemaal rechtlijnig zijn zoals in de eerste poging hierboven. De punten tussen $K$ en $M$ ($M$ zelf niet meegerekend) liggen dichter bij rand $d$ dan bij punt $F$. Het lijnstuk $[KM]$ moet vervangen worden door punten even ver van punt $F$ als van rand $d$. Dit wordt dus een stuk parabool met brandpunt $F$ en richtlijn $d$. Analoog moet $[KL]$ vervangen worden door een stuk parabool met brandpunt $F$ en richtlijn $e$. Verder moet het lijnstuk $[KF]$ weggelaten worden. Dit is geen kamlijn. Zandkorrels die boven $[KF]$ liggen, twijfelen niet waar ze naartoe moeten rollen. Ze rollen naar punt $F$, het dichtste punt van de rand.}

%     \emph{Op de volgende pagina vind je de juiste oplossing.}
        
%    \begin{multicols}{2}
       
%     \includegraphics[angle=-90, width=6.8cm]{UW4201/img/Zand 31.JPG}{}{}
%     \afbeelding[6.8cm]{UW4201/img/Zand 32.JPG}{}{}
    
%    \end{multicols}
            

%     \emph{Als je met hoogtelijnen werkt, zie je dat de zandhoop bij het snijpunt van de twee paraboolstukken, inderdaad iets hoger ligt dan de rechte kamlijnen.}

% \end{vraagenantwoord}     
% \end{lesactiviteit}

% \begin{lesactiviteit}{Zand op een ellips}

% \begin{vraagenantwoord}
%     \vraag{Giet zoveel mogelijk droog zand op een ellips die in karton uitgeknipt is en die horizontaal op een bekertje rust. Wat zie je?}

%     \antwoord{In bovenaanzicht zien we een kamlijn: een stuk van de grote as, symmetrisch rond het middelpunt. Deze kamlijn lijkt niet tot aan de toppen op de grote as te gaan. Een eerste hypothese: misschien vallen de eindpunten van (het bovenaanzicht van) de kamlijn samen met de brandpunten van de ellips? In vooraanzicht zien we dat deze kam niet horizontaal loopt, maar in het midden hoger is dan opzij. Dit is niet zo verwonderlijk: de punten in het midden van de ellips liggen verder van de rand...}
    
% \begin{multicols}{2}

%         \afbeelding[6cm]{UW4201/img/Zand 35.jpg}{}{}
    
%         \afbeelding[6cm]{UW4201/img/Zand 36.jpg}{}{}
        
% \end{multicols}

  

%     \end{vraagenantwoord}

% Vanaf nu werken we in het bovenaanzicht. De punten van de kamlijn beschouwen we als punten van de grote as van de ellips. We kunnen wel verwijzen naar de afrollende zandkorrels in de ruimte, maar als we punten of rechten benoemen, doen we aan vlakke meetkunde.


%     \begin{vraagenantwoord}[resume]
 

%     \vraag{Net als in de vorige opgaven kun je de kamlijn bekijken als een meetkundige plaats van punten die aan een bepaalde voorwaarde voldoen.  Aan welke voorwaarde moeten de punten van de kamlijn voldoen?}

%     \antwoord{De kamlijn is de meetkundige plaats van de punten op gelijke loodrechte afstand tot twee of meer randen van de figuur, hier dus tot de ellips. Immers, zandkorrels rollen naar de dichtste rand. Zandkorrels aan de ene kant van de kamlijn rollen naar de ene rand en zandkorrels aan de andere kant van de kamlijn rollen naar de andere rand.}
%     \emph{Op de volgende figuur klopt dit duidelijk voor de punten van de grote as $[DE]$ die niet te ver van het middelpunt liggen. Punten van de grote as die ver aan de zijkant liggen, dicht bij een van de toppen $D$ of $E$, die rollen recht naar deze top. Ze liggen dus niet op een kamlijn want er is maar één punt van de rand waar ze het dichtste bij liggen. De vraag zal dus zijn: waar zijn de eindpunten $A$ en $B$ van de kamlijn, dus zo dat een punt $P$ tussen $A$ en $B$ twee dichtste randpunten $M$ en $M'$ heeft, en een punt tussen $A$ en $D$ (of tussen $B$ en $E$) slechts één, namelijk $D$ (of $E$). }

%      \afbeelding[6cm]{UW4201/img/Zand 37}{}{}

%     \vraag{Je kunt het misschien omkeren. In plaats van een variabel punt $P$ op de grote as te nemen, en je af te vragen hoe dicht het bij de top moet liggen opdat deze top het dichtste randpunt zou zijn, kun je ook een variabel punt $M$ van de ellips nemen, en dat steeds dichter naar de top $D$ laten gaan, en bepalen waar het bijbehorende punt $P$ van de grote as in de limiet naartoe gaat. Onderzoek dit met GeoGebra voor de ellips $\frac{x^2}{25}+\frac{y^2}{9}=1$.} 

%     \antwoord{Leerlingen tekenen deze ellips met GeoGebra, nemen een punt $M$ hierop, en met de knoppen `raaklijnen' en `loodlijn' kunnen ze de normaal in $M$ tekenen en snijden met de grote as. Ze vinden dat voor $M$ heel dicht bij de top $D$ het punt $P$ coördinaat $(3,2; 0)$ heeft (zie figuur rechts). Dit is nog geen exacte methode, maar het lijkt er alvast op dat het eindpunt $A$ van de kam niet het brandpunt zal zijn. Deze ellips heeft brandpunten $F_{1, 2} (\pm4, 0)$. }


%     \begin{multicols}{2}
    
%          \afbeelding[6cm]{UW4201/img/Zand 38}{}{}
         
%          \afbeelding[6cm]{UW4201/img/Zand 39}{}{}
             
%     \end{multicols}

%     \vraag{We willen nu het eindpunt $A$ aan de rechterkant van de kam exact localiseren. Ken je een eigenschap van de ellips waar de normaal een rol in speelt?}
%     \antwoord{De weerkaatsingseigenschap zegt dat $F_1 M$ en $F_2 M$ gelijke hoeken vormen met de normaal $MP$. (Waardoor een lichtstraal die in $F_1$ vertrekt, in $F_2$ toekomt en omgekeerd...). Dus: $MP$ is de bissectrice van de hoek $\hat{F_1 M F_2}$.}

%    \afbeelding[8.5cm]{UW4201/img/Zand 40}{}{}
       
%     \vraag{Wat zegt dit over de positie van punt $P$ op de grote as. Of, preciezer gezegd, wat weet je over de verhouding $\frac{|F_1 P|}{|F_2 P|}$?}

%     \antwoord{De bissectricestelling zegt: de bissectrice van een hoek van een driehoek verdeelt de overstaande zijde in stukken die in dezelfde verhouding staan als de aanliggende zijden. Hier, in de driehoek $F_1 M F_2$ geeft dit:
%     $$\frac{|F_1 P|}{|F_2 P|}=\frac{|F_1 M|}{|F_2 M|}.$$}
%     \vraag{Maak de redenering af: waar komt $P$ terecht als $M$ naar de top $D$ gaat? }

%     \antwoord{Als $M$ naar $D$ gaat, komt $P$ in het grenspunt $A$ van de kamlijn $[AB]$ terecht. Het punt $A$ is dus bepaald door
%     $$\frac{|F_1 A|}{|F_2 A|}=\frac{|F_1 D|}{|F_2 D|}.$$}
%     \vraag{Maak het eens concreet, voor de ellips van vraag 3.}

%     \antwoord{Voor de ellips $\frac{x^2}{25}+\frac{y^2}{9}=1$ is de halve grote as $5$ en de halve kleine as $3$. We hebben $D (5, 0)$, $F_1(4, 0)$ en $F_2(-4, 0)$. Stel nu $A(u, 0)$. We zoeken $u$ met de gevonden gelijkheid van verhoudingen.
%     $$\frac{4-u}{u-(-4)}=\frac{5-4}{5-(-4)}$$}
%     \emph{Deze vergelijking oplossen naar $u$ geeft: $u=3,2$. Precies wat we met GeoGebra vonden.
%     Merk op: in de vergelijking hierboven gingen we ervan uit dat $u<4$. Dit hebben de leerlingen overtuigend gezien bij het onderzoeken met GeoGebra.}

%     \end{vraagenantwoord}

% \end{lesactiviteit}

% % TODO: bronnen nog deftig integreren?
% \begin{bronnen}
% \item Dendoncker, J.-M. (2025). Moet er nog zand zijn? \emph{Wiskunde en Onderwijs, 51/203}, 28-34.
% \item Élèves des collèges Vieux colombier (LeMans) et Vieux Chêne. (2009). La géométrie du tas de sable. \emph{\url{http://soleil.rene.free.fr/neo/Butinages/La_geometrie_du_tas_de_sable.pdf}}
% \item Eggermont, H., Roelens, M. (2004). Kegelsneden. \emph{Uitwiskeling 20/4}, 9-42.
% \item Gilbert, T. (2017). Tas de sable sur une ellipse. In Gilbert, T., Ninove, L., le GEM, \emph{Le plaisir de chercher en mathématiques}. Louvain-la-Neuve: Presses Universitaires de Louvain.
% \item Levrie, P., Penne, R. (2015). Kegelsneden made easy. \emph{Uitwiskeling 31/3}, 15-17.
% \item Matemateca. (z.d.). Sandpiles. \emph{\url{https://matemateca.ime.usp.br/acervo/montanhas_areia_ingles.html}}
% \item M3 Matemática Multimídia. (23 maart 2012). Montanhas geométricas. \emph{YouTube op \url{https://www.youtube.com/watch?v=f2kBia-jNfg}}
% \item Pierson, P. (2012). La géométrie des tas de sable. \emph{Travail de fin d’études}, Haute École Galilée.
% \end{bronnen}

\end{document}